% 原书 PDF 151，印刷页 139。PDF 152（印刷页 140）为空白。
\part[局部性质与点的 Hilbert 概形]{局部性质与点的 Hilbert 概形\\[24pt]{\large\normalfont\itshape Barbara Fantechi 与 Lothar Göttsche}}
% 原书 PDF 153，印刷页 141。PDF 154（印刷页 142）为空白。
\chapter*{引言}
\addcontentsline{toc}{chapter}{引言}
第六章初步介绍形变理论，随后用它研究 Hilbert 概形与 Quot 概形的局部性质，例如切空间、维数和非奇异性。其思想是，研究从局部 Artin 代数的谱到概形 $X$、并将闭点映到 $p$ 的映射，就能得到关于 $X$ 在点 $p$ 处局部性质的大量信息。当概形 $X$ 自然地表示某个函子时，这一方法尤为有用；Hilbert 概形就是一个例子。这时，需要研究由点 $[Z]\in X$ 所参数化的对象 $Z$ 的无穷小形变，即局部 Artin 代数的谱上的族，其在闭点上的限制是 $Z$。事实表明，切空间与障碍理论能够包含大量信息：需要考察函子的切空间（即 $\Spec k[\varepsilon]/(\varepsilon)^2$ 上的族），以及将形变从一个局部 Artin 代数扩张到稍大的局部 Artin 代数时的障碍。利用切空间与障碍理论，至少可以描述 $X$ 中点 $[Z]$ 的一个形式邻域。

第七章研究点的 Hilbert 概形 $X^{[n]}$，它参数化光滑拟射影簇 $X$ 上长度为 $n$ 的有限子概形。当 $X$ 是曲面时，由于点的 Hilbert 概形具有良好的性质，并与数学的许多其他领域，包括层的模空间、顶点代数，乃至理论物理，都有联系，近年来受到了广泛关注。

我们将构造从 $n$ 点的 Hilbert 概形到第 $n$ 个对称幂的 Hilbert--Chow 态射，它将一个概形映到带有重数的支集。我们研究 $n$ 较小时的几个简单情形，并证明光滑曲线上点的 Hilbert 概形同构于相应的对称幂。随后，我们研究曲面上点的 Hilbert 概形 $S^{[n]}$ 的一个自然分层，并用它概述 Betti 数的计算。事实表明，这些 Betti 数会自然地组织成一个对所有 $n$ 求和的生成函数。这说明不同 $n$ 对应的 $S^{[n]}$ 的上同调群应当紧密联系在一起。例如，Heisenberg 代数在这些上同调群对所有 $n$ 的直和上的作用，就给出了这种联系。本讲义的最后一节简述这一结果，由此引出 Hilbert 概形与顶点代数的联系。

感谢 Leila Khatami 和 Angelo Vistoli 在交换代数方面给予的帮助。
